
@@M0@@ and @@M1@@ are conjugate; in other words, there is an invertible matrix @@M2@@ @@M3@@ such that
@@M4@@
Indeed, the matrix @@M5@@ with @@M6@@, @@M7@@, and otherwise @@M8@@, answers the question (recall that @@M9@@ and @@M10@@ are invertible by hypothesis and @@M11@@). Thus, for any homomorphism @@M12@@ from @@M13@@ to a sheaf of abelian groups @@M14@@, one has
@@M15@@
If @@M16@@ is invertible, this and (6.3.1) imply that @@M17@@ for every @@M18@@, so that @@M19@@ is a product of commutators. If the residue field @@M20@@ of @@M21@@ is not the prime field @@M22@@, then for every invertible @@M23@@ one can find a Zariski neighborhood @@M24@@ of @@M25@@ and an invertible element @@M26@@ such that @@M27@@ is invertible. It follows that @@M28@@ is a product of commutators on @@M29@@.

If @@M30@@, then @@M31@@ is an étale neighborhood of @@M32@@ whose residue field differs from @@M33@@. The result follows.

\sgaStableTarget{sga6.e10.sec.6.5}\subsection*{6.5. Remarks}
\sgaStableTarget{sga6.e10.para.6-5-1}\textbf{6.5.1.} When none of the residue fields of @@M34@@ is isomorphic to @@M35@@, the preceding proof shows that 6.2.2. remains valid for the ringed topos of \emph{Zariski} sheaves on @@M36@@. On the other hand, already for @@M37@@, @@M38@@ differs from the commutator group of @@M39@@. The order of the latter is @@M40@@, whereas the order of @@M41@@ is @@M42@@ (cf. [1]).

\sgaStableTarget{sga6.e10.para.6-5-2}\textbf{6.5.2.} On an arbitrary locally ringed topos @@M43@@, it is not true in general that every representation @@M44@@ is defined by a
